% !TeX root = ../main.tex
\section{Related Work}
\paragraph{Computational histories of research fields.}
\citet{hall2008} used topic models to follow ideas through computational linguistics. Venue retrospectives have also used co-word analysis at CHI \citep{liu2014} and traced technologies, methods, and values at CSCW \citep{wallace2017}. These studies show that a field's history depends on what evidence the retrospective follows.

\paragraph{Reflection within social computing.}
The FAccT retrospective combined publication analysis with a community questionnaire \citep{facct2022}. \citet{weird2024} examined 494 ICWSM papers from 2018--2022, retaining 420 for geographical analysis. They report that 37\% focused exclusively on Western populations, compared with 76\% at CHI and 84\% at FAccT. They also find that ICWSM studies still predominantly examine populations from countries that are more educated, industrialized, and rich than those studied at FAccT. We complement this population-focused analysis with a longer view of topics and problem framings.

\paragraph{Computational social science and its evidence.}
Foundational accounts describe the opportunities and institutional conditions for studying social behavior through digital traces \citep{lazer2009,lazer2020}. Others emphasize platform selection, data access, and the difficulty of treating observed users as a population \citep{ruths2014,p14517}. \citet{wallach2018} distinguishes social-scientific explanation from applying computational methods to social data. These concerns motivate our attention to the questions pursued within a topic.

\paragraph{Measurement, text, and causal inference.}
Topic-model fit does not establish that people find a component meaningful \citep{chang2009}, and automated text analysis requires substantive interpretation and validation \citep{grimmer2013}. Measurement theory further distinguishes an intended construct from its observable operationalization \citep{jacobs2021}. Dictionary estimates can also diverge from survey measures across populations \citep{jaidka2020estimating}. Text-based causal research faces additional problems when language represents confounders, treatments, or outcomes \citep{keith2020}. We draw on these distinctions when comparing framing cues with research designs. Our separation of topics and problems is also informed by problematization, which examines how scholars question assumptions within a body of work \citep{alvesson2011}. Framing cues approximate these questions; they cannot recover the full process through which a problem is formulated.
